\section{Code Generators}
\Java code that is automatically generated by another piece of software should implement this coding conventions in the same way as \bdq{hand crafted} code. Most generators are highly configurable and/or allow to use code templates that are customisable.\footnote{Perhaps you want to have a look to \autocite{TQUADRAT_ORG_FOUNDATION_JAVACOMPOSER}.}

Especially if base classes will be generated it is important that proper comments are generated at least for all API elements; this covers all \lstinline|public| and \lstinline|protected| elements, and in some cases all package-local elements, too.

Java~6 introduced an annotation \lstinline|javax.annotation.Generated|\index{javax.annotation!Generated} (\lstinline|@Generated|) that should be used to mark generated code. Refer to \autocite{ORACLE_DOC_GENERATED_ANNOTATION} for the details.\footnote{In case you still have to use Java~5 as the source code version, you should consider to create your own \lstinline|@Generated| annotation along the lines given by that existing implementation in Java~6 and later.}

%==============================================================================
% End of file!!
%==============================================================================
